\epigraph{The Theory of Interest, as Determined by Impatience to Spend Income and Opportunity to Invest It.}{\textit{Irving Fisher}, title (1930)}

% ==============================================================
\section{The Two-Period Problem}
\label{sec:two-period}
% ==============================================================

We now apply the generalized means framework to intertemporal choice. The key insight is that consumption at different dates is analogous to consumption of different goods---time replaces variety. Interest rates play the role of relative prices, and the tools developed for static CES optimization carry over directly.

% --------------------------------------------------------------
\subsection{From Varieties to Time}
\label{subsec:varieties-to-time}
% --------------------------------------------------------------

In previous sections, an agent chose quantities $\bx = (x_1, \ldots, x_n)$ of different goods subject to a budget constraint $\sum_i p_i x_i = W$. The relative price $p_i/p_j$ governed the trade-off between goods $i$ and $j$.

We now consider an agent who lives for two periods, $t \in \{0, 1\}$. The agent chooses consumption $(c_0, c_1)$ over time, rather than across varieties. The fundamental trade-off is no longer between two goods (e.g., apples and oranges), but between consuming today versus tomorrow, which can be thought of as two goods.

\paragraph{The Interest Rate as a Relative Price.}
If the agent saves one unit of the consumption good today, it grows to $R$ units tomorrow, where $R = 1 + r$ is the gross interest rate. Equivalently, one unit of future consumption costs $q \equiv R^{-1}$ units of current consumption. 

The price $q$ is the relative price of future consumption in terms of current consumption. This price is the bridge between static and dynamic optimization: with $p_0 = 1$ (current consumption as numeraire) and $p_1 = q = R^{-1}$, the intertemporal problem is a special case of the canonical CES problem with two goods.

% --------------------------------------------------------------
\subsection{The Sequential Formulation}
\label{subsec:sequential}
% --------------------------------------------------------------

The standard macroeconomic formulation presents the problem in sequential form.

\paragraph{Preferences.}
The agent has preferences with discount factor $\beta \in (0,1)$ and period utility power utility, $u(c) = \frac{c^{1-1/\theta} - 1}{1 - 1/\theta}$, where $\theta > 0$ is the intertemporal elasticity of substitution (IES). The role if the -1 terms subtracting is to guarantee that utility is continuous in the parameter $\theta$: without this term $U$ is discountinuos at $\theta=0.$ Of course, adding a constant does not change the preference relationship of the agent.

The objective is:
\begin{equation}
U = u(c_0) + \beta \, u(c_1) = \frac{c_0^{1-1/\theta} - 1}{1 - 1/\theta} + \beta \frac{c_1^{1-1/\theta} - 1}{1 - 1/\theta}.
\label{eq:objective-sequential}
\end{equation}

\paragraph{Budget Constraints.}
The agent enters period $0$ with initial wealth $a_0$ and receives income $y_0$. After consuming $c_0$, the remainder is saved in a bond $b_1$:
\begin{equation}
c_0 + b_1 = a_0 + y_0.
\label{eq:budget-t0}
\end{equation}
In period $1$, the bond matures with gross return $R$, and the agent receives income $y_1$:
\begin{equation}
c_1 = R b_1 + y_1.
\label{eq:budget-t1}
\end{equation}
We assume no bequest motive, so the agent consumes all available resources in period $1$. 

We can perform a change of variables and denote by $a_1=R b_1$ the wealth carried to period 1. In that case, the first period's budget is,
\[
c_0 + \frac{a_1}{R} = a_0 + y_0.
\]
Both formulations are equivalent. 

\paragraph{The Optimization Problem.}
The agent maximizes \cref{eq:objective-sequential} subject to the sequential budgets \cref{eq:budget-t0,eq:budget-t1}.

% --------------------------------------------------------------
\subsection{The Lifetime Budget}
\label{subsec:lifetime-budget}
% --------------------------------------------------------------

The two sequential constraints can be consolidated into a single lifetime budget constraint. From \cref{eq:budget-t1}, $b_1 = (c_1 - y_1)/R$. Substituting into \cref{eq:budget-t0}:
\[
c_0 + \frac{c_1 - y_1}{R} = a_0 + y_0 \quad \Longrightarrow \quad c_0 + \frac{c_1}{R} = a_0 + y_0 + \frac{y_1}{R}.
\]
This relationship tells us the present value of consumption equals the present value of wealth. 

\begin{definition}[Lifetime Wealth]
\label{def:lifetime-wealth}
The agent's lifetime wealth is the present discounted value of all resources:
\begin{equation}
W \equiv a_0 + y_0 + \frac{y_1}{R} = a_0 + y_0 + q \cdot y_1.
\label{eq:lifetime-wealth}
\end{equation}
\end{definition}

The lifetime budget constraint becomes:
\begin{equation}
c_0 + q \cdot c_1 = W, \qquad q \equiv R^{-1}.
\label{eq:lifetime-budget}
\end{equation}
This is a linear constraint in $(c_0, c_1)$ with prices $(1, q)$ and wealth $W$---exactly the structure of the static problems.

\begin{calloutnote}[The Role of $q$]
The bond price $q = R^{-1}$ is the relative price of future consumption. When $R$ rises, $q$ falls: future consumption becomes cheaper in terms of current consumption. This is the price effect that drives intertemporal substitution.
\end{calloutnote}

\begin{figure}[H]
  \centering
  \includegraphics[width=0.85\textwidth]{two_period_budget_optimum}
  \caption{Lifetime budget constraint $c_0 + c_1/R = W$, indifference curve, endowment, and optimum, for $\beta = 0.96$, $\theta = 1$, $R = 1.04$.}
  \label{fig:two-period-budget}
\end{figure}

% --------------------------------------------------------------
\subsection{Transformation to Generalized Mean Form}
\label{subsec:transformation-canonical}
% --------------------------------------------------------------

The objective \cref{eq:objective-sequential} is not yet in a generalized mean form. However, we now apply monotone transformations to rewrite it as a CES aggregator, enabling direct application of the solution formulas from \cref{sec:problem-solution}. We present both the classical and canonical representations, showing explicitly how the discount factor $\beta$ maps to aggregator weights.

\paragraph{Step 1: Dropping Constants.}
Since the $-1/(1-1/\theta)$ terms do not depend on $(c_0, c_1)$, we can forget about these additive constants. We check later whether the transformed problems are continuous in $\theta$, which is why the constants were introduced in the first place. 

Maximizing $U$ is equivalent to maximizing:
\begin{equation}
\tilde{U} = \frac{c_0^{1-1/\theta}}{1 - 1/\theta} + \beta \frac{c_1^{1-1/\theta}}{1 - 1/\theta} = \frac{1}{1-1/\theta} \left( c_0^{1-1/\theta} + \beta \, c_1^{1-1/\theta} \right).
\label{eq:U-tilde}
\end{equation}

\paragraph{Step 2: Identifying the Power.}
Define $\rho \equiv 1 - 1/\theta$. The relationship between the power $\rho$ and the IES $\theta$ mirrors the relationship between the power and the elasticity of substitution in the static problem:
\begin{equation}
\rho = 1 - \frac{1}{\theta}, \qquad \theta = \frac{1}{1 - \rho}.
\label{eq:rho-theta}
\end{equation}

The objective becomes:
\begin{equation}
\tilde{U} = \frac{1}{\rho} \left( c_0^{\rho} + \beta \, c_1^{\rho} \right).
\label{eq:U-power}
\end{equation}
The term in parentheses, $c_0^{\rho} + \beta \, c_1^{\rho}$, is almost the inner sum of a generalized mean---but the coefficients $(1, \beta)$ do not sum to one, and they do not yet have the form required by either the classical or canonical representation.

\paragraph{Step 3: Normalizing to  obtain a Classical Form.}
To obtain a classical generalized mean, we need weights $(\tilde{\omega}_0, \tilde{\omega}_1)$ that sum to one. Multiply $\tilde{U}$ by the positive constant $(1+\beta)^{-1}$:
\begin{equation}
\frac{\tilde{U}}{1+\beta} = \frac{1}{\rho} \left( \frac{1}{1+\beta} \, c_0^{\rho} + \frac{\beta}{1+\beta} \, c_1^{\rho} \right) = \frac{1}{\rho} \left( \tilde{\omega}_0 \, c_0^{\rho} + \tilde{\omega}_1 \, c_1^{\rho} \right),
\label{eq:U-classical-inner}
\end{equation}
where:
\begin{equation}
\tilde{\omega}_0 \equiv \frac{1}{1+\beta}, \qquad \tilde{\omega}_1 \equiv \frac{\beta}{1+\beta}, \qquad \tilde{\omega}_0 + \tilde{\omega}_1 = 1.
\label{eq:classical-weights-def}
\end{equation}

\paragraph{Step 4: Applying the Outer Power.}
The inner sum $\bar{\Ssum} = \tilde{\omega}_0 \, c_0^{\rho} + \tilde{\omega}_1 \, c_1^{\rho}>0$ defines the classical generalized mean via $\Mbar = \bar{\Ssum}^{1/\rho}$. 

However, before we raise $\frac{\tilde{U}}{1+\beta}$ to the power $1/\rho$, we must make sure this is a monotone transformation. 

If $\rho>0$, we are fine because $\frac{1}{\rho}s_1>\frac{1}{\rho}s_2\Rightarrow \left(\frac{1}{\rho}s_1\right)^{1/\rho}>\left(\frac{1}{\rho}s_2\right)^{1/\rho}$. When $\rho<0$, we $s_1>s_2\Rightarrow m^{1/\rho}<n^{1/\rho}$, so raising to the power is not a monotone transformation. However, multiplying by -1 and raising to the power $1/\rho$ is a monotone transformation. 
\[
s_1>s_2 \Rightarrow -\frac{1}{\rho}s_1>-\frac{1}{\rho}s_2\Rightarrow \left(-\frac{1}{\rho}s_1\right)^{1/\rho}>\left(-\frac{1}{\rho}s_2\right)^{1/\rho}
\]. 
Since $\left(\frac{1}{\rho}\right)^{1/\rho}$ and $\left(-\frac{1}{\rho}\right)^{1/\rho}$ are positive constants when $\rho$ is positive and negative, respectively, we can drop the constants, again monotone transformations. Hence, maximizing $\tilde{U}$ is equivalent to maximizing the classical mean $\Mbar$.

We have further seen that maximizing a problem in its classical form is equivalent to maximizing it in its canonical form.

\paragraph{The Two Representations.}
We can now state the intertemporal problem in both classical and canonical form. The two are equivalent representations of the same optimization problem, connected by the weight transformation from \cref{sec:problem-solution}.

\medskip
\noindent
\begin{minipage}[t]{0.48\textwidth}
\centering\textbf{Classical Form}
\begin{equation}
\max_{c_0, c_1} \; \Mbar(c_0, c_1; \tilde{\bom}, \rho)
\label{eq:classical-obj-inter}
\end{equation}
where
\begin{equation}
\Mbar = \left( \tilde{\omega}_0 \, c_0^{\rho} + \tilde{\omega}_1 \, c_1^{\rho} \right)^{1/\rho}
\label{eq:classical-mean-inter}
\end{equation}
with weights
\begin{align}
\tilde{\omega}_0 &= \frac{1}{1+\beta} \label{eq:classical-w0} \\[4pt]
\tilde{\omega}_1 &= \frac{\beta}{1+\beta} \label{eq:classical-w1}
\end{align}
\end{minipage}%
\hfill
\begin{minipage}[t]{0.48\textwidth}
\centering\textbf{Canonical Form}
\begin{equation}
\max_{c_0, c_1} \; \M(c_0, c_1; \bom, \rho)
\label{eq:canonical-obj-inter}
\end{equation}
where
\begin{equation}
\M = \left( \omega_0^{1-\rho} c_0^{\rho} + \omega_1^{1-\rho} c_1^{\rho} \right)^{1/\rho}
\label{eq:canonical-mean-inter}
\end{equation}
with weights
\begin{align}
\omega_0 &= \frac{1}{1+\beta^{\theta}} \label{eq:canonical-w0} \\[4pt]
\omega_1 &= \frac{\beta^{\theta}}{1+\beta^{\theta}} \label{eq:canonical-w1}
\end{align}
\end{minipage}

\medskip
\noindent
Both are subject to the same lifetime budget constraint:
\begin{equation}
c_0 + q \, c_1 = W, \qquad q = R^{-1}.
\label{eq:budget-both}
\end{equation}

\paragraph{Deriving the Canonical Weights.}
The canonical weights $(\omega_0, \omega_1)$ are obtained from the classical weights via the transformation formula established in \cref{subsec:equivalence}:
\begin{equation}
\omega_i = \frac{\tilde{\omega}_i^{1/(1-\rho)}}{\sum_{j} \tilde{\omega}_j^{1/(1-\rho)}}.
\label{eq:weight-transform-inter}
\end{equation}

Since $1/(1-\rho) = \theta$, we have:
\begin{equation}
\tilde{\omega}_0^{\theta} = \left( \frac{1}{1+\beta} \right)^{\theta}, \qquad \tilde{\omega}_1^{\theta} = \left( \frac{\beta}{1+\beta} \right)^{\theta} = \frac{\beta^{\theta}}{(1+\beta)^{\theta}}.
\label{eq:weight-transform-calc}
\end{equation}

The sum is:
\begin{equation}
\tilde{\omega}_0^{\theta} + \tilde{\omega}_1^{\theta} = \frac{1 + \beta^{\theta}}{(1+\beta)^{\theta}}.
\label{eq:weight-sum}
\end{equation}

Therefore:
\begin{equation}
\omega_0 = \frac{\tilde{\omega}_0^{\theta}}{\tilde{\omega}_0^{\theta} + \tilde{\omega}_1^{\theta}} = \frac{1}{1 + \beta^{\theta}}, \qquad \omega_1 = \frac{\tilde{\omega}_1^{\theta}}{\tilde{\omega}_0^{\theta} + \tilde{\omega}_1^{\theta}} = \frac{\beta^{\theta}}{1 + \beta^{\theta}}.
\label{eq:canonical-weights-final}
\end{equation}

\begin{calloutnote}[The Role of $\theta$ in the Weight Transformation]
The canonical weight on future consumption is $\omega_1 = \beta^{\theta}/(1+\beta^{\theta})$, not simply $\beta/(1+\beta)$. The IES $\theta$ governs how the discount factor translates into aggregator weights:
\begin{itemize}[nosep]
\item When $\theta = 1$ (log utility), $\omega_1 = \beta/(1+\beta)$: classical and canonical weights coincide.
\item When $\theta > 1$, the canonical weight $\omega_1$ is more sensitive to $\beta$ than the classical weight.
\item When $\theta < 1$, the canonical weight $\omega_1$ is less sensitive to $\beta$.
\end{itemize}
This is the intertemporal analog of the weight transformation between representations in the static problem.
\end{calloutnote}

\paragraph{Verification: Recovering the Inner Sums.}
We can verify the equivalence by checking that the two inner sums are proportional. The canonical inner sum is:
\begin{equation}
\Ssum = \omega_0^{1-\rho} c_0^{\rho} + \omega_1^{1-\rho} c_1^{\rho}.
\label{eq:canonical-inner-inter}
\end{equation}

Using $1 - \rho = 1/\theta$ and the canonical weights:
\begin{align}
\omega_0^{1-\rho} &= \left( \frac{1}{1+\beta^{\theta}} \right)^{1/\theta}, \\[4pt]
\omega_1^{1-\rho} &= \left( \frac{\beta^{\theta}}{1+\beta^{\theta}} \right)^{1/\theta} = \frac{\beta}{(1+\beta^{\theta})^{1/\theta}}.
\label{eq:canonical-powers}
\end{align}

The ratio is:
\begin{equation}
\frac{\omega_1^{1-\rho}}{\omega_0^{1-\rho}} = \beta,
\label{eq:weight-ratio}
\end{equation}
which confirms that the canonical inner sum has terms proportional to $c_0^{\rho}$ and $\beta c_1^{\rho}$, matching the original objective up to a constant.

We further verify that as $\rho \rightarrow 0$, our cannonical form becomes:
\[\mathcal{M}=\left(c_0c_1^{\beta}\right)^{\frac{1}{1+\beta}}.
\]
Taking logs, another monotone transformation, leads to the same objective as the original formulation, up to a constant. Both problems are equivalent also at their limit. 

\begin{calloutimportant}[The Mapping to the Canonical Problem]
The intertemporal problem maps to the canonical CES problem with:
\begin{center}
\renewcommand{\arraystretch}{1.4}
\begin{tabular}{@{}lll@{}}
\toprule
\textbf{Static Problem} & \textbf{Intertemporal Problem} & \textbf{Formula} \\
\midrule
Goods $\{1, \ldots, n\}$ & Periods $\{0, 1\}$ & --- \\
Quantities $x_i$ & Consumption $c_t$ & --- \\
Prices $p_i$ & $(p_0, p_1) = (1, q)$ & $q = R^{-1}$ \\
Wealth $W$ & Lifetime wealth $W$ & $W = a_0 + y_0 + q \cdot y_1$ \\
Elasticity $\varepsilon$ & IES $\theta$ & --- \\
Power $\rho$ & Power $\rho$ & $\rho = 1 - 1/\theta$ \\
Classical weight $\tilde{\omega}_i$ & Normalized discount & $\tilde{\omega}_1 = \beta/(1+\beta)$ \\
Canonical weight $\omega_i$ & Transformed weight & $\omega_1 = \beta^{\theta}/(1+\beta^{\theta})$ \\
\bottomrule
\end{tabular}
\end{center}
With this mapping, all results from \cref{sec:problem-solution}---optimal allocations, budget shares, price indices---apply directly to the intertemporal problem.
\end{calloutimportant}

% --------------------------------------------------------------
\subsection{The Solution}
\label{subsec:two-period-solution}
% --------------------------------------------------------------

With the problem in canonical form, we apply the solution formulas from \cref{sec:problem-solution}.

\paragraph{The Intertemporal Price Index.}
Using that $\frac{\rho}{\rho-1}=\frac{1-\frac{1}{\theta}}{-\frac{1}{\theta}}=1-\theta$, we reconstruct the ideal price index.  The ideal price index aggregates prices $(1, q)$ with weights $(\omega_0, \omega_1)$:
\begin{equation}
\Pstar = \left( \omega_0 \cdot 1^{1-\theta} + \omega_1 \cdot q^{1-\theta} \right)^{1/(1-\theta)} = \left( \omega_0 + \omega_1 \, q^{1-\theta} \right)^{1/(1-\theta)}.
\label{eq:price-index-intertemporal}
\end{equation}
Since $q = R^{-1}$, we have $q^{1-\theta} = R^{\theta - 1}$, and:
\begin{equation}
\Pstar = \left( \omega_0 + \omega_1 \, R^{\theta - 1} \right)^{1/(1-\theta)}.
\label{eq:price-index-R}
\end{equation}
The term $\Pstar$ in this case indicates the cost of transforming a unit of lifetime wealth into lifetime utility (in its canonical form). 

\paragraph{Budget Shares.}
Recall that the optimal expenditure shares were given by:
\[
w_i=\omega_i\left(\frac{p_i}{P^{*}}\right)^{\frac{\rho}{\rho-1}}.
\]
Thus, the expenditure shares on current and future consumption are, consumption at the different dates is given by:
\begin{align}
\wexp_0 &= \frac{c_0}{W} = \omega_0 \left( \frac{1}{\Pstar} \right)^{1-\theta} = \frac{\omega_0}{\omega_0 + \omega_1 \, R^{\theta-1}}, \label{eq:share-c0} \\[6pt]
\wexp_1 &= \frac{q \cdot c_1}{W} = \omega_1 \left( \frac{q}{\Pstar} \right)^{1-\theta} = \frac{\omega_1 \, R^{\theta-1}}{\omega_0 + \omega_1 \, R^{\theta-1}}. \label{eq:share-c1}
\end{align}

\paragraph{Optimal Consumption.}
Applying the demand formula $x_i^* = \omega_i (p_i / \Pstar)^{-\theta} \cdot W / \Pstar$:
\begin{align}
c_0^* &= \omega_0 \left( \frac{1}{\Pstar} \right)^{-\theta} \cdot \frac{W}{\Pstar} = \frac{\omega_0}{(\Pstar)^{1-\theta}} \cdot W = \wexp_0 \cdot W, \label{eq:c0-star} \\[6pt]
c_1^* &= \omega_1 \left( \frac{q}{\Pstar} \right)^{-\theta} \cdot \frac{W}{\Pstar} = \omega_1 \, R^{\theta} \cdot \frac{W}{(\Pstar)^{1-\theta} \cdot \Pstar} = \frac{\wexp_1}{q} \cdot W = \wexp_1 \cdot R \cdot W. \label{eq:c1-star}
\end{align}

\paragraph{The Consumption Ratio.}
Dividing $c_1^*$ by $c_0^*$:
\begin{equation}
\frac{c_1^*}{c_0^*} = \frac{\omega_1}{\omega_0} \left( \frac{q}{1} \right)^{-\theta} = \frac{\omega_1}{\omega_0} R^{\theta}.
\label{eq:consumption-ratio}
\end{equation}
Using the relationship $\omega_1 / \omega_0 = \beta^{1/(1-\rho)} = \beta^{\theta}$:
\begin{equation}
\frac{c_1^*}{c_0^*} = (\beta R)^{\theta}.
\label{eq:consumption-ratio-final}
\end{equation}

This is the fundamental relationship governing consumption growth: when $\beta R > 1$, the agent prefers to defer consumption; when $\beta R < 1$, the agent prefers to consume early.

% --------------------------------------------------------------
\subsection{The Euler Equation}
\label{subsec:euler}
% --------------------------------------------------------------

The Euler equation is the intertemporal analog of the static optimality condition (MRS equals price ratio).

\paragraph{Derivation from the Ratio Condition.}
The FOC ratio condition states:
\begin{equation}
\frac{\omega_0^{1-\rho} c_0^{\rho-1}}{\omega_1^{1-\rho} c_1^{\rho-1}} = \frac{p_0}{p_1} = \frac{1}{q} = R.
\label{eq:ratio-intertemporal}
\end{equation}

Rearranging:
\begin{equation}
\frac{\omega_0^{1-\rho}}{\omega_1^{1-\rho}} \left( \frac{c_0}{c_1} \right)^{\rho - 1} = R.
\label{eq:ratio-rearranged}
\end{equation}

Using $\omega_1^{1-\rho} / \omega_0^{1-\rho} = \beta$ and $\rho - 1 = -1/\theta$:
\begin{equation}
\frac{1}{\beta} \left( \frac{c_0}{c_1} \right)^{-1/\theta} = R \quad \Longrightarrow \quad \left( \frac{c_1}{c_0} \right)^{-1/\theta} = \beta R.
\label{eq:euler-derivation}
\end{equation}

\paragraph{The Euler Equation.}
Recognizing that $u'(c) = c^{-1/\theta}$ for our utility specification:
\begin{equation}
u'(c_0) = \beta R \, u'(c_1) \qquad \Longleftrightarrow \qquad c_0^{-1/\theta} = \beta R \, c_1^{-1/\theta}.
\label{eq:euler}
\end{equation}

\begin{calloutnote}[Interpretation]
The Euler equation states that the marginal utility of consuming one unit today equals the marginal utility of saving that unit, earning return $R$, and consuming tomorrow, discounted by $\beta$. At the optimum, the agent is indifferent between these two options.

Equivalently, the intertemporal marginal rate of substitution equals the gross interest rate:
\[
\text{MRS}_{0,1} = \frac{u'(c_0)}{\beta \, u'(c_1)} = R.
\]
\end{calloutnote}

\begin{figure}[H]
  \centering
  \includegraphics[width=0.85\textwidth]{two_period_euler_growth}
  \caption{Log consumption growth $\log(c_1^\star / c_0^\star)$ as a function of $\log(\beta R)$, for different IES values. The slope equals $\theta$.}
  \label{fig:two-period-euler-growth}
\end{figure}

% --------------------------------------------------------------
\subsection{Special Cases and the Role of the IES}
\label{subsec:special-cases}
% --------------------------------------------------------------

Different values of the IES $\theta$ yield qualitatively different responses to interest rate changes. The Cobb-Douglas case ($\theta = 1$) serves as a crucial benchmark.

\paragraph{Cobb-Douglas ($\rho=0,\theta = 1$, Log Utility).}
When $\theta = 1$, the utility function is $u(c) = \log c$. The consumption ratio becomes:
\begin{equation}
\frac{c_1}{c_0} = \beta R,
\label{eq:ratio-log}
\end{equation}
and the budget shares are constant:
\begin{equation}
\wexp_0 = \frac{1}{1+\beta}, \qquad \wexp_1 = \frac{\beta}{1+\beta}.
\label{eq:shares-log}
\end{equation}

The expenditure share on current consumption is independent of $R$. This is a striking result: changes in the interest rate have no effect on how much wealth the agent allocates to period-$0$ consumption.

\paragraph{Why the Cobb-Douglas Share is Invariant.}
This result connects directly to the cross-price elasticity analysis from the static CES problem. Recall that in the canonical problem, the expenditure share on good $i$ is:
\begin{equation}
\wexp_i = \omega_i \left( \frac{p_i}{\Pstar} \right)^{1-\varepsilon}.
\label{eq:share-recall}
\end{equation}

The elasticity of the expenditure share with respect to another good's price $p_j$ depends on $(1 - \varepsilon)$. When $\varepsilon = 1$ (Cobb-Douglas), we have $1 - \varepsilon = 0$, and the expenditure share becomes:
\begin{equation}
\wexp_i = \omega_i \quad \text{(constant, independent of prices)}.
\label{eq:share-cobb-douglas}
\end{equation}

In the intertemporal context, with $\varepsilon = \theta$ and the price of future consumption $p_1 = q = R^{-1}$, the expenditure on current consumption $c_0$ has zero cross-price elasticity with respect to $q$. The intertemporal expenditure share $\wexp_0$ is invariant to $R$ at $\theta = 1$, mirroring the static Cobb--Douglas result.

To see this in the language of the Slutsky decomposition we develop in \cref{subsec:slutsky}, the Cobb--Douglas knife-edge is exactly where the compensated substitution $-\theta\wexp_1$ and the purchasing-power income effect $+\wexp_1$ cancel each other---these are the only two forces operating in the static problem with exogenous wealth.

\medskip
\noindent
\emph{But $c_0$ is not invariant in $R$.} The intertemporal problem has a third channel that the static problem lacks: \emph{endowment revaluation}. When $R$ rises, the present value of $y_1$ falls, and lifetime wealth $W$ shrinks. Since $c_0 = \wexp_0 \cdot W$ with $\wexp_0$ constant at $\theta = 1$, the Cobb-Douglas response to $R$ is fully channeled through $W$ rather than through the relative price. This is the simplest illustration of the point that the intertemporal problem is an endowment economy: a price change revalues the agent's holdings of future income, and this revaluation does not disappear at the Cobb-Douglas benchmark even though the substitution and purchasing-power channels do.

\begin{calloutimportant}[The Cobb-Douglas Benchmark]
The case $\theta = 1$ is the knife-edge where:
\begin{itemize}[nosep]
\item Expenditure shares $\wexp_0, \wexp_1$ are constant across all values of $R$.
\item Cross-price elasticities of expenditure are zero.
\item In the Slutsky decomposition, the compensated substitution and the purchasing-power income channels exactly offset each other.
\end{itemize}
This is not a coincidence---it is the same mathematical structure as in the static CES problem with unit elasticity. What survives at $\theta = 1$ in the intertemporal problem, and what is absent from the static problem, is the endowment-revaluation channel: changes in $R$ continue to move $c_0$ through their effect on lifetime wealth $W$. The two problems share identical \emph{share} comparative statics at the Cobb--Douglas benchmark, but the intertemporal level response carries an extra wealth term that the static problem does not.
\end{calloutimportant}

\paragraph{Low IES ($\theta < 1$).}
When $\theta < 1$, the cross-price elasticity of the expenditure share is positive: $1 - \theta > 0$. A fall in $q$ (rise in $R$) \emph{increases} the expenditure share on current consumption. The income effect dominates the substitution effect. Intuitively, the agent is relatively unwilling to substitute across time; when the return to saving rises, the agent feels wealthier and consumes more today.

\paragraph{High IES ($\theta > 1$).}
When $\theta > 1$, the cross-price elasticity is negative: $1 - \theta < 0$. A fall in $q$ \emph{decreases} the expenditure share on current consumption. The substitution effect dominates. The agent readily substitutes toward the now-cheaper future consumption, saving more today.

\paragraph{Perfect Substitutes ($\theta \to \infty$).}
As $\theta \to \infty$, utility becomes linear in consumption, and the agent is willing to substitute perfectly across time. The substitution effect completely dominates: if $\beta R > 1$, the agent saves everything; if $\beta R < 1$, the agent borrows to the limit.

\begin{table}[H]
\centering
\caption{Consumption Ratios for Different Utility Specifications}
\label{tab:utility-cases}
\renewcommand{\arraystretch}{1.6}
\begin{tabular}{@{}lll@{}}
\toprule
\textbf{Utility} & \textbf{Euler Equation} & \textbf{Consumption Ratio} \\
\midrule
CRRA: $\dfrac{c^{1-1/\theta}}{1-1/\theta}$ & $c_0^{-1/\theta} = \beta R \, c_1^{-1/\theta}$ & $\dfrac{c_1}{c_0} = (\beta R)^{\theta}$ \\[12pt]
Log: $\log c$ & $\dfrac{1}{c_0} = \beta R \dfrac{1}{c_1}$ & $\dfrac{c_1}{c_0} = \beta R$ \\[12pt]
Exponential: $-\sigma e^{-c/\sigma}$ & $e^{(c_1 - c_0)/\sigma} = \beta R$ & $c_1 - c_0 = \sigma \log(\beta R)$ \\[8pt]
\bottomrule
\end{tabular}
\end{table}

\begin{calloutnote}[Beyond CES: Other Utility Classes]
The CRRA family (including log utility as a limiting case) and exponential utility share a common structure: they admit closed-form Euler equations with clean comparative statics. Other specifications, such as linear-quadratic utility, also yield tractable Euler equations but do not fit the CES/generalized mean framework. These alternative classes have their own uses---linear-quadratic preferences, for instance, lead to certainty equivalence in stochastic settings---but lie outside the scope of these notes.
\end{calloutnote}

\begin{calloutnote}[Consumption Smoothing]
In all CES cases, if $\beta R = 1$, the agent chooses $c_0 = c_1$: perfect consumption smoothing. This benchmark is useful for understanding deviations---when $\beta R \neq 1$, consumption tilts toward the period with higher effective return.
\end{calloutnote}

\begin{figure}[H]
  \centering
  \includegraphics[width=0.85\textwidth]{two_period_c0_vs_R_by_theta}
  \caption{Current consumption $c_0^\star$ as a function of $R$ for different values of $\theta$, with $y_0 = y_1 = 1$, $\beta = 0.96$.}
  \label{fig:two-period-c0-by-theta}
\end{figure}

% --------------------------------------------------------------
\subsection{Comparative Statics: The Slutsky Decomposition}
\label{subsec:slutsky}
% --------------------------------------------------------------

We now develop a systematic approach to comparative statics with respect to the interest rate. The goal is to decompose the total effect of a change in $R$ into substitution and wealth effects, paralleling the Slutsky decomposition from static consumer theory. Because $R$ enters both the budget constraint (through the relative price $q = 1/R$) and the present value of future income (through $y_1/R$), the decomposition has a richer structure than its static analog.

\begin{figure}[H]
  \centering
  \includegraphics[width=0.85\textwidth]{two_period_wealth_vs_R}
  \caption{Lifetime wealth $W = a_0 + y_0 + y_1/R$ as a function of $R$, for three income profiles. A rate change is a pure wealth shock whose sign depends on the timing of income.}
  \label{fig:two-period-wealth-vs-R}
\end{figure}

\paragraph{The System in Levels.}
The agent's optimum is characterized by the Euler equation and the budget constraint:
\begin{align}
u'(c_0) &= \beta R \, u'(c_1), \label{eq:euler-system} \\
c_0 + \frac{c_1}{R} &= W. \label{eq:budget-system}
\end{align}
Lifetime wealth $W$ itself depends on $R$ through the present value of future income:
\begin{equation}
W = a_0 + y_0 + \frac{y_1}{R}.
\label{eq:wealth-R}
\end{equation}
This is a system of two equations in two unknowns $(c_0, c_1)$, with $R$ as a parameter. Changes in $R$ affect the system through three channels: the intertemporal price in the Euler equation, the relative price in the budget constraint, and the present value of future income.

\paragraph{Differentiating the System.}
We work in elasticity form, expressing all changes as log-deviations. Differentiating the Euler equation:
\begin{equation}
\frac{u''(c_0)}{u'(c_0)} c_0 \cdot \frac{dc_0}{c_0} = \frac{u''(c_1)}{u'(c_1)} c_1 \cdot \frac{dc_1}{c_1} + \frac{dR}{R}.
\label{eq:euler-diff}
\end{equation}
Define the coefficient of relative risk aversion (the inverse of the IES) as
\begin{equation}
\sigma(c) \equiv -\frac{u''(c) \cdot c}{u'(c)}.
\label{eq:rra-def}
\end{equation}
For CRRA utility $\sigma(c) = 1/\theta$ is constant. The differentiated Euler equation becomes
\begin{equation}
-\sigma_0 \cdot \frac{dc_0}{c_0} = -\sigma_1 \cdot \frac{dc_1}{c_1} + \frac{dR}{R},
\label{eq:euler-diff-sigma}
\end{equation}
where $\sigma_0 = \sigma(c_0)$ and $\sigma_1 = \sigma(c_1)$.

Differentiating the budget constraint:
\begin{equation}
c_0 \cdot \frac{dc_0}{c_0} + \frac{c_1}{R} \cdot \frac{dc_1}{c_1} - \frac{c_1}{R} \cdot \frac{dR}{R} = dW.
\label{eq:budget-diff}
\end{equation}
Differentiating lifetime wealth:
\begin{equation}
dW = -\frac{y_1}{R} \cdot \frac{dR}{R}.
\label{eq:wealth-diff}
\end{equation}
Combining:
\begin{equation}
c_0 \cdot \frac{dc_0}{c_0} + \frac{c_1}{R} \cdot \frac{dc_1}{c_1} = \frac{c_1 - y_1}{R} \cdot \frac{dR}{R}.
\label{eq:budget-combined}
\end{equation}

\paragraph{Matrix Formulation and Algebraic Channel Split.}
Collecting terms, the system can be written in matrix form:
\begin{equation}
\underbrace{\begin{bmatrix} -\sigma_0 & \sigma_1 \\[4pt] c_0 & \dfrac{c_1}{R} \end{bmatrix}}_{\mathbf{A}}
\begin{bmatrix} \dfrac{dc_0}{c_0} \\[8pt] \dfrac{dc_1}{c_1} \end{bmatrix}
= \begin{bmatrix} 1 \\[4pt] \dfrac{c_1 - y_1}{R} \end{bmatrix} \frac{dR}{R}.
\label{eq:matrix-system}
\end{equation}
The right-hand side admits the algebraic split:
\begin{equation}
\begin{bmatrix} 1 \\[4pt] \dfrac{c_1 - y_1}{R} \end{bmatrix}
= \underbrace{\begin{bmatrix} 1 \\[4pt] \dfrac{c_1}{R} \end{bmatrix}}_{\substack{\text{relative-price}\\\text{vector at fixed }W}}
\;+\; \underbrace{\begin{bmatrix} 0 \\[4pt] -\dfrac{y_1}{R} \end{bmatrix}}_{\substack{\text{endowment}\\\text{revaluation}}}.
\label{eq:channel-split}
\end{equation}
The first vector is the impact of $R$ on the Euler equation and the relative price in the budget, holding lifetime wealth $W$ \emph{nominally} fixed. The second vector captures the erosion of present-value wealth, $-y_1/R$. Applying $\mathbf{A}^{-1}$ linearly, the total response decomposes into these two algebraic pieces.

\medskip
\noindent
\emph{A note on terminology.} It is tempting to call the first vector the ``substitution effect'' and the second the ``income effect''---and in much of the macro literature this is what is done. We will see momentarily that this labeling is a slight abuse: the response at \emph{fixed nominal $W$} is not the same as the response at \emph{fixed utility}, which is the proper Hicksian definition. The two coincide only at the Cobb--Douglas benchmark $\theta = 1$. We will return to this point after computing both pieces.

\paragraph{The CRRA Case.}
With CRRA utility, $\sigma_0 = \sigma_1 = \sigma = 1/\theta$. The matrix $\mathbf{A}$ and its inverse simplify considerably:
\begin{equation}
\mathbf{A} = \begin{bmatrix} -\sigma & \sigma \\[4pt] c_0 & \dfrac{c_1}{R} \end{bmatrix}, \qquad
\det(\mathbf{A}) = -\sigma \left( c_0 + \frac{c_1}{R} \right) = -\sigma W,
\label{eq:A-crra}
\end{equation}
and
\begin{equation}
\mathbf{A}^{-1} = \frac{1}{W} \begin{bmatrix} -\dfrac{1}{\sigma} \dfrac{c_1}{R} & 1 \\[4pt] \dfrac{1}{\sigma} c_0 & 1 \end{bmatrix}.
\label{eq:A-inverse}
\end{equation}

\paragraph{Fixed-$W$ Response.}
Applying $\mathbf{A}^{-1}$ to the relative-price vector:
\begin{equation}
\mathbf{A}^{-1} \begin{bmatrix} 1 \\[4pt] \dfrac{c_1}{R} \end{bmatrix}
= \frac{1}{W} \begin{bmatrix} \dfrac{c_1}{R}\left(1 - \dfrac{1}{\sigma}\right) \\[8pt] \dfrac{c_0}{\sigma} + \dfrac{c_1}{R} \end{bmatrix}.
\label{eq:sub-effect-calc}
\end{equation}
Using $1 - 1/\sigma = 1 - \theta$ and the expenditure shares $\wexp_0 = c_0/W$, $\wexp_1 = c_1/(RW)$:
\begin{equation}
\left.\begin{bmatrix} \dfrac{dc_0}{c_0} \\[8pt] \dfrac{dc_1}{c_1} \end{bmatrix}\right|_{\text{fixed-}W} = \begin{bmatrix} (1-\theta) \, \wexp_1 \\[4pt] \theta \, \wexp_0 + \wexp_1 \end{bmatrix} \frac{dR}{R}.
\label{eq:fixed-W-effect}
\end{equation}
This is the response that would obtain if lifetime nominal wealth $W$ were held constant during the price change. For $c_0$, the fixed-$W$ response is $(1-\theta)\wexp_1$: negative when $\theta>1$, positive when $\theta<1$, and zero at $\theta=1$.

\paragraph{Endowment Revaluation.}
Applying $\mathbf{A}^{-1}$ to the endowment-revaluation vector:
\begin{equation}
\mathbf{A}^{-1} \begin{bmatrix} 0 \\[4pt] -\dfrac{y_1}{R} \end{bmatrix}
= -\frac{y_1}{RW} \begin{bmatrix} 1 \\[4pt] 1 \end{bmatrix}.
\label{eq:wealth-effect-calc}
\end{equation}
Both consumption levels adjust proportionally---a pure income response whose magnitude is the PV-share of future income in lifetime wealth, $y_1/(RW)$. The effect is non-positive whenever $y_1 > 0$: a higher interest rate lowers the present value of future income, regardless of preferences.
\begin{equation}
\left.\begin{bmatrix} \dfrac{dc_0}{c_0} \\[8pt] \dfrac{dc_1}{c_1} \end{bmatrix}\right|_{\text{reval.}} = -\frac{y_1}{RW} \begin{bmatrix} 1 \\[4pt] 1 \end{bmatrix} \frac{dR}{R}.
\label{eq:income-effect}
\end{equation}

\paragraph{The Marshallian Total.}
Combining the two algebraic pieces gives the total response:
\begin{equation}
\;\begin{bmatrix} \dfrac{dc_0}{c_0} \\[8pt] \dfrac{dc_1}{c_1} \end{bmatrix} = \underbrace{\begin{bmatrix} (1-\theta) \, \wexp_1 \\[4pt] \theta \, \wexp_0 + \wexp_1 \end{bmatrix}}_{\substack{\text{fixed-$W$}\\\text{response}}} \frac{dR}{R} \;-\; \underbrace{\frac{y_1}{RW} \begin{bmatrix} 1 \\[4pt] 1 \end{bmatrix}}_{\substack{\text{endowment}\\\text{revaluation}}} \frac{dR}{R}.\;
\label{eq:slutsky-complete}
\end{equation}
For current consumption alone:
\begin{equation}
\frac{d\log c_0}{d\log R} = (1-\theta)\,\wexp_1 \;-\; \frac{y_1}{RW}.
\label{eq:slutsky-c0}
\end{equation}
A useful cross-check: differencing the two rows of \cref{eq:slutsky-complete} gives $d\log c_1/d\log R - d\log c_0/d\log R = \theta\,(\wexp_0 + \wexp_1) = \theta$, exactly what one obtains by differentiating the consumption ratio $c_1/c_0 = (\beta R)^\theta$ directly. The decomposition is therefore consistent with the Euler equation at all points.

% --------------------------------------------------------------
\paragraph{Three Primitive Channels.}
The response in \cref{eq:slutsky-c0} mixes three economically distinct forces, and to make sense of the saver/borrower split we need to disentangle them. A change in $R$ does three things at once:

\begin{enumerate}
\item \textbf{Compensated substitution.} Holding utility constant, the cheaper good ($c_1$, since $q = 1/R$ falls when $R$ rises) gets a larger share of consumption and the dearer good ($c_0$) a smaller share. For CRRA utility, the Hicksian response of $c_0$ to $\log R$ is $-\theta\,\wexp_1$.
\item \textbf{Purchasing-power income effect.} Even at fixed nominal $W$, the agent's existing bundle gets cheaper when $c_1$ becomes cheaper, so real income rises. This is the standard Slutsky income term, $-x_j \,\partial x_i^M / \partial W$ from static consumer theory; in our problem it equals $+\wexp_1$ for $c_0$.
\item \textbf{Endowment revaluation.} Lifetime wealth $W$ is itself a function of $R$, since $dW/dR = -y_1/R^2 < 0$. This is the channel \emph{absent} from the static problem with exogenous $W$: it appears here because the agent owns a portion of their consumption stream. The induced effect on $c_0$ is $-y_1/(RW)$.
\end{enumerate}

The three add up to the total:
\begin{equation}
\frac{d\log c_0}{d\log R} \;=\; \underbrace{-\theta\,\wexp_1}_{\substack{\text{compensated}\\\text{substitution}}} \;+\; \underbrace{\wexp_1}_{\substack{\text{purchasing-power}\\\text{income}}} \;-\; \underbrace{\frac{y_1}{RW}}_{\substack{\text{endowment}\\\text{revaluation}}}.
\label{eq:three-term}
\end{equation}
The arithmetic is immediate: $-\theta\wexp_1 + \wexp_1 = (1-\theta)\wexp_1$, recovering \cref{eq:slutsky-c0}.

\begin{calloutimportant}[Why an Extra Channel?]
In the static consumer problem with exogenous wealth, only two forces operate: the compensated substitution and the purchasing-power income effect. Their sum is the Marshallian response. In an \emph{endowment economy}---and the intertemporal problem is one, with $y_1$ the period-1 endowment---a price change also revalues the agent's holdings. The full Slutsky decomposition therefore needs three terms, and any honest two-term grouping has to specify how the three are combined.
\end{calloutimportant}

% --------------------------------------------------------------
\paragraph{Two Equivalent Two-Term Groupings.}
The three primitive channels can be packaged into two natural Slutsky-style forms, each pairing the compensated substitution with one or both of the income-type terms.

\medskip
\noindent
\textbf{Form A (Marshallian, fixed-$W$ split).} Group the compensated substitution with the purchasing-power income effect, then keep the endowment revaluation as a separate ``wealth'' piece:
\begin{equation}
\frac{d\log c_0}{d\log R} \;=\; \underbrace{(1-\theta)\,\wexp_1}_{\substack{\text{response at}\\\text{fixed nominal }W}} \;-\; \underbrace{\frac{y_1}{RW}}_{\substack{\text{endowment}\\\text{revaluation}}}.
\label{eq:slutsky-fixedW}
\end{equation}
This is \cref{eq:slutsky-c0}. Its first term is the response that obtains if we pretend $W$ is exogenous and apply the two-good Slutsky formula directly. It is \emph{not} the Hicksian compensated effect; the two coincide only at the Cobb--Douglas benchmark $\theta = 1$.

\medskip
\noindent
\textbf{Form B (Hicksian + net-demand income).} Group the two income-type terms together. From $a = (c_1 - y_1)/R$ and $\wexp_1 = c_1/(RW)$ we obtain the identity
\begin{equation}
\frac{a}{W} \;=\; \wexp_1 - \frac{y_1}{RW}.
\label{eq:saving-identity}
\end{equation}
Substituting into \cref{eq:three-term}:
\begin{equation}
\boxed{\;\frac{d\log c_0}{d\log R} \;=\; \underbrace{-\theta\,\wexp_1}_{\substack{\text{compensated}\\\text{substitution}}} \;+\; \underbrace{\frac{a}{W}}_{\substack{\text{net-demand}\\\text{income effect}}}.\;}
\label{eq:slutsky-hicksian}
\end{equation}
This is the proper Hicksian Slutsky decomposition for the intertemporal problem. The first term holds utility constant, exactly as in the static Slutsky equation. The second is the income effect of the price change, scaled by the agent's \emph{net demand} for the good whose price moved: $c_1 - y_1 = R\,a$.

\begin{calloutnote}[Two Forms, One Decomposition]
\cref{eq:slutsky-fixedW} and \cref{eq:slutsky-hicksian} are algebraically identical:
\[
(1-\theta)\,\wexp_1 - \frac{y_1}{RW} \;=\; -\theta\,\wexp_1 + \frac{a}{W}.
\]
They differ only in how the three primitive channels of \cref{eq:three-term} are grouped:
\begin{itemize}[nosep]
\item Form A treats $\wexp_1$ as part of the ``substitution'' piece (giving $(1-\theta)\wexp_1$) and $-y_1/(RW)$ as the wealth effect.
\item Form B treats $\wexp_1$ as part of the income piece (giving $a/W$) and isolates $-\theta\wexp_1$ as the genuine compensated substitution.
\end{itemize}
Form A is the form one obtains by applying the matrix machinery directly. Form B is the one that maps cleanly onto the textbook Slutsky decomposition for an endowment economy, and onto the geometric Slutsky picture in which the indifference curve $U_0$ is held fixed.
\end{calloutnote}

\begin{figure}[H]
  \centering
  \includegraphics[width=\textwidth]{two_period_slutsky_geometric}
  \caption{Geometric Slutsky decomposition in the $(c_0, c_1)$ plane, for a saver (left, $y_0 = 1.6$, $y_1 = 0.4$) and a borrower (right, $y_0 = 0.4$, $y_1 = 1.6$), for a rate increase from $R_0 = 1.04$ to $R_1 = 1.25$ with $\beta = 0.96$, $\theta = 2$. Point $A$ is the old optimum, $B$ the Hicksian bundle on $U_0$ at the new price (compensated substitution, $A \to B$), $C$ the new optimum (combined income effect, $B \to C$). The budget line rotates around the endowment $\omega$, leaving $A$ inside the new budget for the saver and outside for the borrower.}
  \label{fig:two-period-slutsky-geometric}
\end{figure}

\begin{figure}[H]
  \centering
  \includegraphics[width=\textwidth]{two_period_slutsky_hicksian}
  \caption{Form A (Marshallian fixed-$W$ split) of \cref{eq:slutsky-fixedW}: $d\log c_0 / d\log R = (1-\theta)\wexp_1 - y_1/(RW)$ as a function of $R$, for the saver (left) and borrower (right), $\beta = 0.96$, $\theta = 2$. The fixed-$W$ piece is identical across agents; the endowment-revaluation piece differs in magnitude (larger for the borrower, whose $W$ is smaller and whose $y_1$ is larger).}
  \label{fig:two-period-slutsky-hicksian}
\end{figure}

\begin{figure}[H]
  \centering
  \includegraphics[width=\textwidth]{two_period_slutsky_price_savings}
  \caption{Form B (Hicksian + net-demand income) of \cref{eq:slutsky-hicksian}: $d\log c_0 / d\log R = -\theta\,\wexp_1 + a/W$. The compensated substitution channel is identical across agents; the net-demand income channel $a/W$ flips sign between savers and borrowers, making the saver/borrower asymmetry transparent.}
  \label{fig:two-period-slutsky-price-savings}
\end{figure}

\paragraph{Welfare: The Envelope Theorem.}
The consumption response describes behavior. The \emph{welfare} consequences admit a still cleaner expression. Let $V(R) = u(c_0^\star(R)) + \beta\, u(c_1^\star(R))$ denote indirect utility. Applying the envelope theorem to the lifetime budget---only the price of $c_1$ and the endowment $y_1$ are affected by $R$---gives
\begin{equation}
\frac{dV}{dR} = \beta\, u'(c_1^\star) \cdot a.
\label{eq:envelope}
\end{equation}
The sign of $dV/dR$ is the sign of $a$ alone: savers gain from a rate increase, borrowers lose, regardless of the IES. The IES affects the magnitude of the welfare change (through how $u'(c_1^\star)$ and $a$ themselves respond) but not its sign. This is the cleanest possible distillation of the distributional content of an interest-rate change.

\begin{figure}[H]
  \centering
  \includegraphics[width=0.85\textwidth]{two_period_welfare}
  \caption{Welfare $V(R) - V(R_0)$ for a saver and a borrower as $R$ varies. By the envelope theorem, $dV/dR = \beta \, u'(c_1^\star) \cdot a$, which carries the sign of savings. The asymmetry is robust to $\theta$.}
  \label{fig:two-period-welfare}
\end{figure}

\begin{calloutnote}[Who Gains, Who Loses?]
A rate increase makes savers ($a > 0$) better off and borrowers ($a < 0$) worse off, by the envelope formula \cref{eq:envelope}. The consumption response can be read in either of the two equivalent two-term forms:
\begin{itemize}[nosep]
\item \emph{Form A} \cref{eq:slutsky-fixedW}: the fixed-$W$ piece $(1-\theta)\wexp_1$ is common across agents; the endowment-revaluation piece $-y_1/(RW)$ is larger in magnitude for the borrower, whose lifetime wealth is more concentrated in $y_1$.
\item \emph{Form B} \cref{eq:slutsky-hicksian}: the compensated substitution $-\theta\wexp_1$ is common across agents; the net-demand income channel $a/W$ flips sign between savers and borrowers.
\end{itemize}
In both readings, a borrower's $c_0$ falls by more than a saver's. Welfare, by the envelope theorem, follows the sign of $a$ directly.
\end{calloutnote}

\paragraph{Connection to the Static Slutsky Equation.}
Form B \cref{eq:slutsky-hicksian} is the intertemporal instance of the standard Slutsky equation for an \emph{endowment economy}:
\begin{equation}
\frac{\partial x_i^M}{\partial p_j} \;=\; \underbrace{\frac{\partial h_i}{\partial p_j}}_{\substack{\text{Hicksian}\\\text{substitution}}} \;-\; \underbrace{(x_j^\star - e_j)\,\frac{\partial x_i^M}{\partial W}}_{\substack{\text{net-demand}\\\text{income effect}}},
\label{eq:slutsky-static}
\end{equation}
where $h_i$ is the Hicksian (compensated) demand and $e_j$ is the endowment of good $j$. In the textbook consumer problem $e_j = 0$ and the income term reduces to $-x_j^\star \partial x_i^M / \partial W$, scaled by gross consumption. In an endowment economy, the agent receives endowment $e_j$ valued at $p_j$, so a price change revalues both the consumption bundle and the endowment. The relevant net position is the difference, $x_j^\star - e_j$.

In our intertemporal problem, $j = 1$ refers to future consumption, $p_1 = q = 1/R$, and $e_1 = y_1$. The net demand is $c_1 - y_1 = R\,a$. The savings position $a$ is therefore the natural intertemporal counterpart of static net demand, and its sign is the sign of whether the agent is a net buyer or a net seller of future consumption. A saver (net buyer of $c_1$) is hurt by a price increase on $c_1$---which is exactly a fall in $R$. A borrower (net seller of $c_1$) gains. \cref{eq:slutsky-hicksian} states this precisely.

\paragraph{Application: Capital Income Taxation.}
The Slutsky decomposition clarifies the effects of taxing capital income---a central question in public finance. Suppose the government imposes a proportional tax $\tau$ on capital income, so the after-tax gross return is $R(1-\tau)$ rather than $R$. From the agent's perspective, this is equivalent to a reduction in the interest rate.

Consider two policy designs:

\medskip
\noindent
\textit{Tax without rebate.} The government collects revenue $\tau R a$ (where $a = a_0 + y_0 - c_0$ is savings) and does not return it. The agent faces the after-tax return $R(1-\tau)$ and the full Slutsky response applies: both the fixed-$W$ response and the endowment revaluation operate.

\medskip
\noindent
\textit{Tax with lump-sum rebate.} The government collects the same revenue but returns it as a lump-sum transfer that exactly offsets the change in lifetime wealth, so $W$ is unchanged. Only the fixed-$W$ response operates.

\medskip
This distinction maps onto fixed-$W$ versus uncompensated responses:

\begin{center}
\renewcommand{\arraystretch}{1.4}
\begin{tabular}{@{}lll@{}}
\toprule
\textbf{Policy} & \textbf{Effect on $c_0$} & \textbf{Response} \\
\midrule
Tax with $W$-preserving rebate & $(1-\theta) \, \wexp_1 \cdot \dfrac{d\tau}{1-\tau}$ & Fixed-$W$ (Form A) \\[8pt]
Tax without rebate & Full Slutsky formula & Uncompensated (Marshallian) \\
\bottomrule
\end{tabular}
\end{center}
The literal Hicksian deadweight-loss calculation requires a transfer that keeps \emph{utility} (not $W$) constant, which differs from the lump-sum scheme above except at $\theta = 1$. We treat the $W$-preserving rebate here because it is the natural macroeconomic implementation: revenue is rebated lump-sum, period.

\begin{calloutimportant}[Efficiency versus Distribution]
The welfare cost of capital taxation depends on the \emph{compensated substitution channel} $-\theta\,\wexp_1$ in Form B \cref{eq:slutsky-hicksian}---the genuine intertemporal-allocation distortion caused by the tax wedge. This is governed by the IES $\theta$:
\begin{itemize}[nosep]
\item If $\theta$ is large, the compensated substitution is large in magnitude, and capital taxes are highly distortionary.
\item If $\theta$ is small, agents are unwilling to substitute across time, and capital taxes cause little allocative distortion.
\end{itemize}
The \emph{distributional} impact, by contrast, works through the net-demand income channel of \cref{eq:slutsky-hicksian} and the envelope welfare formula \cref{eq:envelope}: a capital income tax transfers resources from savers to the government (or to borrowers, if the revenue is redistributed). Whether this is desirable depends on the correlation between saving and welfare-relevant characteristics.

The Cobb--Douglas case ($\theta = 1$) is again a knife-edge: the fixed-$W$ response on $c_0$ is zero, so a $W$-preserving capital tax has no effect on current consumption. But the Hicksian compensated channel $-\theta\wexp_1 = -\wexp_1$ is still nonzero---the tax still causes deadweight loss, even at the Cobb--Douglas benchmark. The two ``substitution effects'' diverge here.
\end{calloutimportant}

% --------------------------------------------------------------
\subsection{Summary}
\label{subsec:two-period-summary}
% --------------------------------------------------------------

\begin{table}[H]
\centering
\caption{The Two-Period Consumption-Savings Problem}
\label{tab:two-period-summary}
\renewcommand{\arraystretch}{1.6}
\begin{tabular}{@{}ll@{}}
\toprule
\textbf{Object} & \textbf{Formula} \\
\midrule
Lifetime wealth & $W = a_0 + y_0 + R^{-1} y_1$ \\[6pt]
Lifetime budget & $c_0 + R^{-1} c_1 = W$ \\[6pt]
Euler equation & $u'(c_0) = \beta R \, u'(c_1)$ \\[6pt]
Consumption ratio (CRRA) & $c_1 / c_0 = (\beta R)^{\theta}$ \\[6pt]
Price index & $\Pstar = \left( \omega_0 + \omega_1 R^{\theta-1} \right)^{1/(1-\theta)}$ \\[6pt]
Optimal $c_0$ & $c_0^* = \wexp_0 \cdot W$ \\[6pt]
Optimal $c_1$ & $c_1^* = \wexp_1 \cdot R \cdot W$ \\[6pt]
Savings & $a = a_0 + y_0 - c_0^* = a_0 + y_0 - \wexp_0 \cdot W$ \\[6pt]
\bottomrule
\end{tabular}
\end{table}

\begin{callouttip}[Key Insight]
The two-period consumption-savings problem is a CES optimization problem in disguise:
\begin{itemize}[nosep]
\item Time periods $\{0, 1\}$ play the role of goods $\{1, 2\}$
\item The bond price $q = R^{-1}$ is the relative price of future consumption
\item The IES $\theta$ is the intertemporal elasticity of substitution
\item The discount factor $\beta$ determines the relative weight on future utility
\end{itemize}
All the machinery of generalized means---price indices, budget shares, comparative statics---applies directly.
\end{callouttip}
